%%%%%%%%%%%%%%%%%%%%%%%%%%%%%%%%%%%%%%%%%%%%%%%%%%%%%%%%%%%%%%%%%%%%%%%%
% Beamer Presentation - LaTeX - Template Version 1.0 (10/11/12)
% This template has been downloaded from: http://www.LaTeXTemplates.com
% License: % CC BY-NC-SA 3.0 (http://creativecommons.org/)
% Modified by Rahmat M. Samik-Ibrahim

% REV004 Mon 21 Feb 2022 13:30:00 WIB
% REV003 Mon 14 Feb 2022 21:50:54 WIB
% REV001 Mon 07 Feb 2022 05:54:47 WIB
% STARTX Wed 14 Sep 2016 10:45:19 WIB
%%%%%%%%%%%%%%%%%%%%%%%%%%%%%%%%%%%%%%%%%%%%%%%%%%%%%%%%%%%%%%%%%%%%%%%%%

% PACKAGES AND THEMES 
\documentclass[aspectratio=169, xcolor=table, notheorems, hyperref={pdfpagelabels=false}]{beamer}
\input{beamer.tex}
\newcommand{\revision}{REV021 06-Feb-2023}
% w! tmptmp
% REV021: Mon 06 Feb 2023 00:00
% REV019: Thu 02 Feb 2023 19:00
% REV009: Mon 18 Apr 2022 06:00
% REV007: Thu 17 Mar 2022 10:00
% REV001: Mon 07 Feb 2022 18:00
% STARTS: Wed 24 Aug 2016 19:00
%%%%%%%%%%%%%%%%%%%%%%%%%%%%%%%
\newcommand{\kopikopi}{\textcopyright{}2016-2023 CBKadal + VauLSMorg}



% XXXXXXXXXXXXXXXXXXXXXXXXXXXXXXXXXXXXXXXXXXXXXXXXXXXXXXXXXXXXXXXXXXXXXXXXXX
% The short title appears at the bottom of every slide, 
% the full title is only on the title page
% \date{\today}
\title[\kopikopi]
{CSCE604227 System Programming \\
CSCE604227 Pemrograman Sistem \\
Week 02:
Revisit Linux From Scratch}
\author{C. BinKadal}
\institute[SDN]
{
Sendirian Berhad\\
\medskip
\url{https://sp.vlsm.org/Slides/sp02.pdf}
\\ \texttt{Always check for the latest revision!}
}
\date{\revision}

% XXXXXXXXXXXXXXXXXXXXXXXXXXXXXXXXXXXXXXXXXXXXXXXXXXXXXXXXXXXXXXXXXXXXXXXXXX
\begin{document}

\lstset{
basicstyle=\ttfamily\tiny, % \tiny \small \footnotesize 
breakatwhitespace=true,
language=C,
columns=fullflexible,
keepspaces=true,
breaklines=true,
tabsize=3, 
showstringspaces=false,
extendedchars=true}

\section{Start}
\begin{frame}
\titlepage
\end{frame}

% XXXXXXXXXXXXXXXXXXXXXXXXXXXXXXXXXXXXXXXXXXXXXXXXXXXXXXXXXXXXXXXXXXXXXXXXXX

%%%%%%%%%%%%%%%%%%%%%%%%%%%%%%%
% REV019: Thu 02 Feb 2023 00:00
% REV003: Mon 14 Feb 2022 21:00
% REV001: Mon 07 Feb 2022 18:00
% START0: Wed 14 Sep 2016 10:00
%%%%%%%%%%%%%%%%%%%%%%%%%%%%%%%

\begin{frame}[fragile]
\section{Schedule}
\frametitle{SP221\footnote{%
) This information will be on \textbf{EVERY} page two (2) of this course material.}): 
System Progamming}

\vspace{5pt}

\scalebox{1.2}{%
\begin{tabular}{|c|c|l|}
\hline
\textbf{Week} & 
\textbf{Topic} \\ 
\hline
Week 00  & Overview \\
Week 01  & Linux Kernel and Programming Interface \\
Week 02  & Revisit Linux From Scratch \\
Week 03  & FUSE: Filesystem in Userspace  \\
Week 04  & Project FUSE 1 \\
Week 05  & Project FUSE 2 \\
\hline
Week 06  & Project FUSE 3 \\
Week 07  & Project FUSE 4 \\
Week 08  & Project FUSE 5 \\
Week 09  & Project FUSE 6 \\
Week 10  & Project FUSE Presentation \\
% MidTerm  & 00 XXX 2020 (XX:XX-XX:XX) & MidTerm (UTS) & \cellcolor{red!44} TBA! \\
% Reserved & 00 XXX - 00 XXX 2020 & Q \& A & \\
% Final    & 00 XXX 2020 XX:XX & First Part Final  (UAS tahap I)  & \cellcolor{red!44} This schedule is   \\
% Extra    & NA & No Extra assignment & \cellcolor{red!44} subject to change. \\
\hline
\end{tabular}
}
\end{frame}

\begin{frame}[fragile]
\frametitle{\textbf{STARTING POINT} --- 
{
\definecolor{links}{HTML}{FDEE00}
\hypersetup{colorlinks,linkcolor=,urlcolor=links}
\url{https://sp.vlsm.org/}
}
}
\begin{itemize}
\item[$\square$] \textbf{Text Book} ---
The Linux Programming Interface, 2010, No Starch Press, 
ISBN 978-1-59327-220-3 --- \url{https://man7.org/tlpi/}. 
\item[$\square$] \textbf{Resources}
\begin{itemize}
\item[$\square$] \href{https://scele.cs.ui.ac.id/course/view.php?id=3545}{\textbf{SCELE}} ---
\url{https://scele.cs.ui.ac.id/course/view.php?id=3545}.\\
The enrollment key is \textbf{XXX}.
\item[$\square$] \textbf{Download Slides and Demos from GitHub.com} \\
\url{https://github.com/os2xx/sysprog/}:

                 {\scriptsize%
                 \href{https://sp.vlsm.org/Slides/sp00.pdf}{\texttt{sp00.pdf} (W00)},
                 \href{https://sp.vlsm.org/Slides/sp01.pdf}{\texttt{sp01.pdf} (W01)},
                 \href{https://sp.vlsm.org/Slides/sp02.pdf}{\texttt{sp02.pdf} (W02)},
                 \href{https://sp.vlsm.org/Slides/sp03.pdf}{\texttt{sp03.pdf} (W03)},

                 \href{https://sp.vlsm.org/Slides/sp04.pdf}{\texttt{sp04.pdf} (W04)},
                 \href{https://sp.vlsm.org/Slides/sp05.pdf}{\texttt{sp05.pdf} (W05)},
                 \href{https://sp.vlsm.org/Slides/sp06.pdf}{\texttt{sp06.pdf} (W06)},
                 \href{https://sp.vlsm.org/Slides/sp07.pdf}{\texttt{sp07.pdf} (W07)},

                 \href{https://sp.vlsm.org/Slides/sp08.pdf}{\texttt{sp08.pdf} (W08)},
                 \href{https://sp.vlsm.org/Slides/sp09.pdf}{\texttt{sp09.pdf} (W09)},
                 \href{https://sp.vlsm.org/Slides/sp10.pdf}{\texttt{sp10.pdf} (W10)}.
                 }
\item[$\square$] \textbf{LFS} --- \url{http://www.linuxfromscratch.org/lfs/view/stable/}
\item[$\square$] \textbf{OSP4DISS} --- \url{https://osp4diss.vlsm.org/}
\item[$\square$] \textbf{This is How Me DO IT!} --- \url{https://doit.vlsm.org/}
\begin{itemize}
\item[$\square$] PS: ''Me'' rhymes better than ''I'' duh!
\end{itemize}
\end{itemize}
\end{itemize}
\end{frame}



% XXXXXXXXXXXXXXXXXXXXXXXXXXXXXXXXXXXXXXXXXXXXXXXXXXXXXXXXXXXXXXXXXXXXXXXXXX
% Throughout your presentation, if you choose to use \section{} and
% \subsection{} commands, these will automatically be printed on
% this slide as an overview of your presentation
\section{Agenda}
\begin{frame}{Outline}
  \frametitle{Agenda}
  \tableofcontents[sections={1-}]
\end{frame}
% \begin{frame}
%    \frametitle{Agenda (2)}
%    \tableofcontents[sections={12-}]
% \end{frame}

% XXXXXXXXXXXXXXXXXXXXXXXXXXXXXXXXXXXXXXXXXXXXXXXXXXXXXXXXXXXXXXXXXXXXXXXXXX
\section{LFS: Linux From Scratch}
\begin{frame}
\frametitle{LFS: Linux From Scratch}
\begin{itemize}
\item \href{https://youtu.be/jEoM3qan9Gs}{THIS IS HOW WE DOIT!}
\item \url{http://www.linuxfromscratch.org/lfs/view/stable/}
\item To build a GNU/Linux system from scratch (source code).
\item To learn a GNU/Linux system inside out.
\item To use a Virtual Machine.
\item A Chicken and Egg dependency problem:
\begin{itemize}
\item It would be best if you had the tools to build an Operating System.
\item You need an Operating System to build tools.
\item To build a cross-toolchain (compiler and its libraries).
\item To build cross utilities using the cross-toolchain.
\item To build an Operating System in a chroot environment.
\item To do iterations (if necessary).
\end{itemize}
\item How deep would you like to know of a ''real'' Operating System?
\item \textbf{YOU} decide!
\end{itemize}
\end{frame}

% XXXXXXXXXXXXXXXXXXXXXXXXXXXXXXXXXXXXXXXXXXXXXXXXXXXXXXXXXXXXXXXXXXXXXXXXXX
\section{Gnulib --- The GNU Portability Library}
\begin{frame}
\frametitle{Gnulib --- The GNU Portability Library\footnote{%
) Adopted from \url{https://www.gnu.org/software/gnulib/}%
})}

\begin{itemize}
\item a central location for common GNU code
\item intended to be shared among GNU packages at the source level
\item no distribution tarball 
\item See also:
\begin{itemize}
\item GNU Coding Standards --- \url{https://www.gnu.org/prep/standards/}
\item Information for maintainers --- \url{https://www.gnu.org/prep/maintain/}
\item Licenses --- \url{https://www.gnu.org/licenses/}
\item Config --- \url{https://savannah.gnu.org/projects/config/}
\end{itemize}
\end{itemize}

\end{frame}

% XXXXXXXXXXXXXXXXXXXXXXXXXXXXXXXXXXXXXXXXXXXXXXXXXXXXXXXXXXXXXXXXXXXXXXXXXX
\section{Autotools}
\begin{frame}
\frametitle{Autotools\footnote{%
) Adopted from \url{https://www.gnu.org/software/automake/faq/}%
})}

\begin{itemize}
\item Refers to the software packages, consists of programs:
\begin{itemize}
\item Autoconf: helps create portable configure and testsuite scripts.
\begin{itemize}
\item URL: \url{https://www.gnu.org/software/autoconf/}
\item autoreconf
\item autoconf
\item autoheader
\item autoscan
\end{itemize}
\item Automake: helps create portable Makefiles.
\begin{itemize}
\item URL: \url{https://www.gnu.org/software/automake/}
\item aclocal
\item automake
\end{itemize}
\item Libtool: helps create and use shared libraries portably. 
\begin{itemize}
\item URL: \url{https://www.gnu.org/software/libtool/}
\item libtoolize
\end{itemize}
\end{itemize}
\item See also:
\begin{itemize}
\item Autotools Mythbuster --- \url{https://autotools.info/}
\item \href{https://www.gnu.org/savannah-checkouts/gnu/automake/manual/automake.html\#Hello-World}{Small Autotools ''Hello World'' Example}
\end{itemize}
\end{itemize}

\end{frame}

% XXXXXXXXXXXXXXXXXXXXXXXXXXXXXXXXXXXXXXXXXXXXXXXXXXXXXXXXXXXXXXXXXXXXXXXXXX
\section{Small Autotools "Hello World" Example}
\begin{frame}[fragile]
\frametitle{Small Autotools "Hello World" Example (1)\footnote{%
) Adopted from {\tiny \url{https://www.gnu.org/savannah-checkouts/gnu/automake/manual/automake.html\#Hello-World}}%
})}

\begin{itemize}
\item Try this and push it to \url{https://github.com/os2xx/sharesp221/}
\begin{itemize}
\item Work inside your \texttt{autotools/} folder.
\end{itemize}

\item \textbf{Filename:} \texttt{src/main.c}
% \begin{lstlisting}[basicstyle=\ttfamily\tiny]
\begin{lstlisting}[basicstyle=\ttfamily\small]

#include <config.h>
#include <stdio.h>

int
main (void)
{
  puts ("Hello World!");
  puts ("I am cbkadal and this is " PACKAGE_STRING ".");
  return 0;
}

\end{lstlisting}
\end{itemize}

\end{frame}

% XXXXXXXXXXXXXXXXXXXXXXXXXXXXXXXXXXXXXXXXXXXXXXXXXXXXXXXXXXXXXXXXXXXXXXXXXX
\begin{frame}[fragile]
\frametitle{Small Autotools "Hello World" Example (2)}

\begin{itemize}
\item \textbf{Filename:} \texttt{README}
% \begin{lstlisting}[basicstyle=\ttfamily\tiny]
\begin{lstlisting}[basicstyle=\ttfamily\small]
This is a demonstration package for GNU Automake.
Type 'info Automake' to read the Automake manual.
Makefile.am and src/Makefile.am contain Automake instructions 
for these two directories.
\end{lstlisting}

\item \textbf{Filename:} \texttt{src/Makefile.am}
\begin{lstlisting}[basicstyle=\ttfamily\small]
bin_PROGRAMS = hello
hello_SOURCES = main.c
\end{lstlisting}

\item \textbf{Filename:} \texttt{Makefile.am}
\begin{lstlisting}[basicstyle=\ttfamily\small]
SUBDIRS = src
dist_doc_DATA = README

\end{lstlisting}
\end{itemize}

\end{frame}

% XXXXXXXXXXXXXXXXXXXXXXXXXXXXXXXXXXXXXXXXXXXXXXXXXXXXXXXXXXXXXXXXXXXXXXXXXX
\begin{frame}[fragile]
\frametitle{Small Autotools "Hello World" Example (3)}

\begin{itemize}
\item \textbf{Filename:} \texttt{configure.ac}
\begin{lstlisting}[basicstyle=\ttfamily\small]
AC_INIT([amhello], [1.0], [bug-automake@gnu.org])
AM_INIT_AUTOMAKE([-Wall -Werror foreign])
AC_PROG_CC
AC_CONFIG_HEADERS([config.h])
AC_CONFIG_FILES([
   Makefile
   src/Makefile
])
AC_OUTPUT

\end{lstlisting}

\item \textbf{RUN:}
\begin{lstlisting}[basicstyle=\ttfamily\small]
autoreconf --install
./configure
make
src/hello

\end{lstlisting}
\end{itemize}

\end{frame}

% XXXXXXXXXXXXXXXXXXXXXXXXXXXXXXXXXXXXXXXXXXXXXXXXXXXXXXXXXXXXXXXXXXXXXXXXXX
\end{document}

